\section{复现指南}
\label{sec:reproduce}

\subsection{代码地址}

本项目全部源代码、工具脚本、伪标签生成器与实验日志已发布在 GitHub：

\begin{center}
\fbox{\parbox{0.86\textwidth}{\centering
\url{https://github.com/jiusheng58/cogeci-reproduction}
}}
\end{center}

仓库不含数据集、BERT 权重与训练权重（体积原因），获取方式见下文。

训练好的 5 折权重（\texttt{checkpoints/E1\_full\_rulelabels/fold\{1..5\}\_best.pt}，
各约 535\,MB，合计约 2.6\,GB）因体积过大未随提交包提供，仍保留在本机。
按下文给出的训练命令重新运行即可得到逐位相同的结果
（固定 \texttt{seed=209}，单卡约 4.2 小时）；如评审需要可直接向作者索取。

\subsection{目录结构}

\begin{lstlisting}[basicstyle=\ttfamily\scriptsize]
cogeci-reproduction/
|-- _repo/                        官方仓库（含本项目补丁）
|   |-- main.py                   训练主循环（+ 权重保存、日志路径修正）
|   |-- parameter.py              超参（+ --grad_ckpt, --cgl_w, --cgl_mode）
|   |-- utils.py                  自 iLIF 补齐
|   |-- model/
|   |   |-- model.py              网络定义（5 处性能补丁 + CGL 参数化 + STE 分支）
|   |   |-- _orig_model.py        原版备份（用于等价性验证）
|   |   `-- bert-base-uncased/    需自行下载
|   |-- data/
|   |   |-- database/exist/*.pkl  自 iLIF 补齐
|   |   `-- examples/llm/         伪标签 CSV（本项目生成）
|   `-- checkpoints/<实验名>/     每折最优权重
|-- tools/
|   |-- event_span.py             子词位置 -> 词级 span 反推
|   |-- pseudo_from_rule.py       规则法伪标签生成
|   |-- label_sentences.py        LLM 打标（含断点续跑、关思考）
|   |-- agree_rule_llm.py         规则 vs LLM 一致性对照
|   |-- agree_by_reltype.py       按 intra/inter 拆分对照
|   |-- make_ablation_table.py   从日志生成消融表
|   |-- validate_csv.py           训练前伪标签终检
|   |-- equiv_patch.py            性能补丁数值等价性验证
|   |-- bench_vram.py             显存峰值实测
|   `-- estimate_tokens.py        打标 token 消耗估算
|-- report/                       LaTeX 报告源文件（本文件所在目录）
|-- run_ablation.sh               消融实验串行队列
`-- logs_*.txt                    各实验完整训练日志
\end{lstlisting}

\subsection{复现步骤}

\paragraph{步骤 1：环境}

\begin{lstlisting}[language=bash]
python -m venv .venv
.venv/Scripts/pip install torch==2.6.0 --index-url https://download.pytorch.org/whl/cu126
.venv/Scripts/pip install torch_geometric==2.5.3 transformers==4.46.3 ^
    scikit-learn tqdm networkx openai
\end{lstlisting}

\noindent 注意 \texttt{transformers} 必须为 4.x（5.x 已移除 \texttt{AdamW}）。

\paragraph{步骤 2：补齐缺失文件}
从 iLIF 仓库取得 \texttt{utils.py} 与 \texttt{data/database/exist/*.pkl}，
下载 \texttt{bert-base-uncased} 至 \texttt{\_repo/model/}。

\paragraph{步骤 3：生成伪标签并终检}

\begin{lstlisting}[language=bash]
cd D:/wangbclass/cogeci
for S in dev train; do for V in strict loose; do
  ./.venv/Scripts/python.exe tools/pseudo_from_rule.py --split $S --variant $V \
      --out_csv _repo/data/examples/llm/rule/${S}_${V}.csv
done; done

./.venv/Scripts/python.exe tools/validate_csv.py \
    --pkl _repo/data/database/exist/EventStoryLine_bert_intra_and_inter_train_features.pkl \
    --csv _repo/data/examples/llm/llama3help_train_sen.csv
\end{lstlisting}

\paragraph{步骤 4：训练}
必须在 \texttt{\_repo/} 目录下运行（\texttt{process\_esc\_csr.py} 使用相对路径，
否则读不到伪标签会静默回退全零）。

\begin{lstlisting}[language=bash]
cd _repo
../.venv/Scripts/python.exe -u main.py --dataset_type EventStoryLine \
  --cache_path ./data/database/exist --model_name ./model/bert-base-uncased \
  --no_of_classes 2 --t_lr 2e-5 --lr 1e-4 --w 0.6 --threshold 2 --epoch 30 \
  --batch_size 1 --beta_intra 0.3 --beta_inter 0.7 --max_iteration 9 \
  --min_iteration 2 --num_heads 4 --mlp_drop 0.4 --seed 209 --gamma 2.0 \
  --grad_ckpt --cgl_w 6.0 --method E1_full_rulelabels
\end{lstlisting}

\noindent 消融实验只需改 \texttt{--cgl\_w}：E0 取 \texttt{0}，并把 \texttt{--method} 改为
\texttt{E0\_no\_cgl}；E3 取 \texttt{--cgl\_mode ste}（权重仍为 6），\texttt{--method} 改为
\texttt{E3\_ste\_termgrad}。再加 \texttt{--fold\_start 1 --fold\_end 1} 可只跑 Fold 1。

\paragraph{步骤 5：一致性对照（可选）}

\begin{lstlisting}[language=bash]
./.venv/Scripts/python.exe tools/agree_rule_llm.py --split dev
./.venv/Scripts/python.exe tools/agree_by_reltype.py --split dev
./.venv/Scripts/python.exe tools/equiv_patch.py      # 等价性验证
\end{lstlisting}

\subsection{硬件与耗时参考}

全部实验在 RTX 4050 Laptop（6\,GB）上完成，单折训练约 50 分钟，
5 折合计约 4.2 小时。若显存 $\ge$ 16\,GB，可去掉 \texttt{--grad\_ckpt} 以获得更快的单步速度，
数值结果不变。
